\medskip

\emph{Pour chacune des affirmations suivantes, indiquer si elle est vraie ou fausse.\\
Chaque réponse doit être justifiée. Une réponse non justifiée ne rapporte aucun point.}

\medskip

\begin{enumerate}
\item Soient $E$ et $F$ les ensembles $E = \{1~;~2~;~3~;~4~;~5~;~6~;~7\}$ et 

$F = \{0~;~1~;~2~;~3~;~ 4~;~5~;~6~;~7~;~8~;~9\}$.

\smallskip

\textbf{Affirmation \no 1 :} Il y a davantage de 3-uplets d'éléments distincts de $E$ que de combinaisons à 4 éléments de $F$.
\end{enumerate}

\begin{minipage}{0.65\linewidth}
\begin{enumerate}[resume]
\item Dans le repère orthonormé ci-contre, on a représenté la fonction carré, notée $f$, ainsi que le carré ABCD de côté 3.

\textbf{Affirmation \no 2 :} La zone hachurée et le carré ABCD ont la même aire.
\end{enumerate}
\end{minipage}\hfill
\begin{minipage}{0.3\linewidth}
\psset{unit=0.5cm,arrowsize=2pt 3}
\begin{pspicture}(-4.5,-1)(4.5,10)
\psaxes[linewidth=1.25pt,labelFontSize=\scriptstyle]{->}(0,0)(-4.5,0)(4.5,10)
\psframe[linecolor=blue,linewidth=1.35pt](3,3)
\pscustom[fillstyle=vlines]{\psplot[plotpoints=2000,linewidth=1.25pt,linecolor=red]{0}{3.}{x dup mul}\psline(3,9)(3,0)(0,0)}
\uput[dl](0,0){\blue \small A}\uput[dr](3,0){\blue \small B}\uput[ur](3,3){\blue \small C}\uput[ul](0,3){\blue \small D}
\psplot[plotpoints=2000,linewidth=1.25pt,linecolor=red]{-3.1}{3.1}{x dup mul}
\end{pspicture}
\end{minipage}

\begin{enumerate}[start=3]
\item On considère l'intégrale $J$ ci-dessous : 

\[J = \displaystyle\int_1^2 x \ln (x)\,\text{d}x.\]

\textbf{Affirmation \no 3 :} Une intégration par parties permet d'obtenir: $J = \dfrac{7}{11}$.

\item Sur $\R$, on considère l'équation différentielle 

\[(E) :\quad  y' = 2y - \e^x.\]

\textbf{Affirmation \no 4 :} La fonction $f$ définie sur $\R$ par $f(x) = \e^x + \e^{2x}$ est solution de l'équation différentielle $(E)$.
\item  Soit $x$ donné dans [0~;~1[. On considère la suite $(u_n)$ définie pour tout entier naturel $n$ par:

\[u_n = (x - 1)\e^n + \cos (n).\]

\textbf{Affirmation \no 5 :} La suite $(u_n)$ diverge vers $-\infty.$
\end{enumerate}
